\documentclass[addpoints,ngerman,12pt,answers]{exam}
\usepackage{babel}
\usepackage[top=2cm,bottom=2.5cm,left=1.5cm,right=1.5cm]{geometry}
\usepackage[utf8]{inputenc}
\usepackage[T1]{fontenc}
\usepackage{booktabs}
\usepackage{graphicx}
\usepackage{csquotes}
\usepackage{paralist}
%\usepackage{wasysym}
%\usepackage[math]{iwona}
\usepackage{setspace}
\usepackage{amsmath} 
\usepackage[math]{iwona}
\usepackage{setspace}
\usepackage{enumitem}
\usepackage{amsmath}
\usepackage{amsfonts}


\pointpoints{Point}{Points}
\bonuspointpoints{Bonuspunkt}{Bonuspunkte}
\renewcommand{\solutiontitle}{\noindent\textbf{Solution:}\enspace}

\chqword{Frage}   
\chpgword{Seite} 
\chpword{Punkte}   
\chbpword{Bonus Points} 
\chsword{Erreicht}   
\chtword{Gesamt}

\checkboxchar{\Square}
\checkedchar{\CheckedBox}

\pagestyle{headandfoot}
\runningheadrule
\firstpageheader{ECON 311}{Week F12}{}
\runningheader{ECON 311}{Week F12}{}
\firstpagefooter{ECON 311}{Week F12}{\thepage\,/\,\numpages}
\runningfooter{ECON 311}{Week F12}{\thepage\,/\,\numpages}
\doublespacing

\begin{document}
	\vspace*{3em}
	
	\makebox[\textwidth]{Name:\enspace\hrulefill}
	
	\vspace*{3em}
	\begin{questions}
		\question Nu Burger Bistro is thinking about acquiring a second stove by borrowing \$1,200 from the bank at a 10\% interest rate. The current setup includes only one stove that can cook 24 burgers per hour, and each burger is sold for \$32.
		\begin{enumerate}[label=(\alph*)]
			\item Considering a depreciation rate of 6\%, should Nu Burger Bistro invest in the new stove? Explain your reasoning.
			\item 
			Now, assume Nu Burger Bistro is considering acquiring additional stoves, with each successive stove being less productive than the previous one. The productivity of the kth stove is given by $\frac{24}{K}$, implying that the 24th stove only cooks one burger per hour. Considering the same prices and depreciation rate as in part $a$, how many stoves should Nu Burger Bistro purchase to maximize its profit?
		\end{enumerate}
	\newpage
	\begin{solution}
		\begin{enumerate}[label=(\alph*)]
			\item For the 2nd stove, it increases the production of burger from 24 units to 48 units. Therefore the MPK in dollar term is $24\times \$32 = \$768$. Using the formula,
			\[\frac{MPK}{p_K} = R+\delta -\frac{\Delta p_K}{p_K}\]
			\[\frac{768}{1200} > 0.10+ 0.06 - 0\]
			As the LHS is greater than the RHS, the Bistro should invest in the new stove
			\item We can sub in $24/K$ as MPK. 
			\[\frac{24\times 32}{1200K} = 0.16 \implies K=4\]
		\end{enumerate}
	\end{solution}		

	\end{questions}
	
%	\begin{center}
%		\combinedgradetable[h][questions]
%	\end{center}
	
\end{document}